\section{A model of operation and future research}
\label{sec:model}
\label{sec:model-projection}

The model records the strategies a manager can operate and the capabilities available for further research. The controller observes a filtered belief, operates
from a verified strategy library, and may initiate positive-duration research
that generates and verifies a candidate before admission. The theorem does not
derive library invariance from finance primitives; it identifies the conclusion
that follows when the primitives pass the explicit factorization checks below.
The retention decision is made once, before the controlled process is evaluated.

\subsection{Hidden state, filtered belief, and operating profiles}

Let \(X\) be a finite hidden market-state set with at least two elements and
let \(B\) be a finite nonempty set of filtered beliefs.  A belief
\(b\in B\) has an exact rational interpretation
\(\mu_b\in\Delta_{\mathbb Q}(X)\); the already-filtered belief evolves under
an exact rational Markov kernel \(P_B\).  The controller observes \(b\), not
the hidden state \(x\in X\).

The finite strategy catalog is \(S\), with distinguished inactive strategy
\(s_0\).  Strategy \(s\) has immutable state-contingent operating payoff
\(u_s:X\to\mathbb Q\) and belief-indexed operating profile
\begin{equation}
  j_s(b):=\sum_{x\in X}\mu_b(x)u_s(x),
  \qquad j_{s_0}(b)=0.
  \label{eq:operational-profile}
\end{equation}
Thus \(j_s(b)\) is the payoff from operating \(s\) now; research costs and
continuation values are separate objects.  The formal projection theorem uses
the resulting exact profile table as primitive data, so it does not claim a
kernel-verified derivation of that table from a hidden-state model.

Partial information matters to the controlled projection through the filtered
belief and its transition law. It is not required by the later tagged-cover
equivalence, which needs only a finite index set of operating-profile rows. The
financial study in Section~\ref{sec:financial-case-studies} uses calendar-year
performance profiles crossed with declared cost and risk settings as these
rows; it does not estimate the hidden-state filter in this subsection.

\subsection{Modules, closure, and the raw retained library}

Let \(M\) be a finite module set.  Each strategy carries an immutable set
\(\mods(s)\subseteq M\), with \(\mods(s_0)=\varnothing\).  A module is a
declared reusable capability tag for later research, not a direct operating
payoff or necessarily a separately owned software artifact.  The
declared closure operator
\[
  \cl:2^M\longrightarrow 2^M
\]
is extensive, monotone, and idempotent.  These properties allow capabilities
to be generated from combinations of raw modules; identity closure is the
special case \(\cl(A)=A\).

A raw retained library is a source subset \(L\subseteq S\) containing
\(s_0\).  Its operating frontier, raw module inventory, and generative closure
are
\begin{align}
  F_L(b)&:=\max_{s\in L}j_s(b),
  &U_L^{\mathrm{raw}}&:=\bigcup_{s\in L}\mods(s),
  \label{eq:frontier-union}\\
  C_L&:=\cl(U_L^{\mathrm{raw}}).
  \label{eq:raw-library-closure}
\end{align}
Finiteness and mandatory inactivity make the maximum attained and imply
\(F_L(b)\geq0\).  The raw decision state is \((b,L)\): continuing earns the
frontier payoff, advances the belief once, and leaves the library unchanged.

An entry may itself be a portfolio strategy, but the operating decision here
selects one catalog entry at a time. Equality of \(F_L\) therefore preserves
the best entry's modeled payoff at each belief; it does not by itself preserve
the risk--return opportunities from combining entries with different joint
returns. An application that reallocates across entries must include the
relevant joint payoff and allocation information in its preservation target.

\subsection{Generation, verification, admission, and timing}

Let \(Q\) be a finite project set.  At project \(q\), belief \(b\), and
closure \(C\), candidate generation follows an exact law
\[
  G(q,b,C)\in\Delta_{\mathbb Q}(\{\mathsf{none}\}\mathbin{\dot\cup}S).
\]
A generated candidate \(s\) passes verification with exact probability
\(\nu(q,b,C,s)\).  The admitted law is derived rather than separately chosen:
\begin{equation}
  \Gamma(q,b,C)(s)=G(q,b,C)(s)\nu(q,b,C,s),
  \label{eq:admitted-law-main}
\end{equation}
and all failed-generation and failed-verification mass is assigned to
\(\mathsf{none}\).  Successful admission adds the verified catalog row,
whereas every other outcome leaves the source unchanged:
\begin{equation}
  L\oplus\mathsf{none}=L,
  \qquad L\oplus s=L\cup\{s\}.
  \label{eq:raw-library-update}
\end{equation}

Each available project has an initiation cost, an operation/suspension flag,
and an integer duration \(d_q\geq1\).  Its exact joint completion law couples
the intervening belief path with the admitted outcome; the two components need
not be independent.  The incumbent library remains in force until completion,
then \eqref{eq:raw-library-update} is applied.  Full probability-table
definitions, null-outcome normalization, path and outcome marginals, and
pushforward identities are assigned to Online Resource~1, Section~S1.

\subsection{The productive library state}

\begin{definition}[Frontier--closure library state]
\label{def:finite-library-state}
The productive state of a retained library is
\begin{equation}
  K_L:=(F_L,C_L).
  \label{eq:closure-compression}
\end{equation}
It is the \emph{library component} of the controlled state.  The filtered
belief remains part of the decision state, and an unfinished project also
requires its identity and remaining duration.
\end{definition}

For \(K=(F,C)\), define the local compressed update
\[
 \operatorname{addK}(K,\mathsf{none})=K,
 \qquad
 \operatorname{addK}(K,s)=
 \left(b\mapsto\max\{F(b),j_s(b)\},
       \cl(C\cup\mods(s))\right).
\]
Extensivity and idempotence of closure give
\begin{equation}
  K_{L\oplus o}=\operatorname{addK}(K_L,o).
  \label{eq:raw-compressed-local-update}
\end{equation}

\begin{assumption}[Finite library-invariant research process]
\label{asm:aor-projection-process}
The sets \(B,S,M,Q\) are finite, \(B\) and every admissible library are
nonempty, and all displayed profiles and probability tables are exact.  The
inactive strategy has zero profile and no modules; closure is extensive,
monotone, and idempotent.  Generation and primitive verification depend on the
incumbent library only through \((q,b,C_L)\).  Project availability, initiation
cost, duration, operation flag, and the joint completion law depend on the
library only through \((b,K_L)\).  Each duration is positive, the joint law has
the declared belief-path and admitted-outcome marginals, and discount satisfies
\(0\leq\beta<1\).  Stationary conclusions additionally assume the recorded
raw and compressed Bellman operators are contractions.
\end{assumption}

\section{What safe retention preserves}
\label{sec:preservation}

\begin{theorem}[Raw-to-compressed controlled Markov projection]
\label{thm:raw-to-compressed-projection}
Under \cref{asm:aor-projection-process}, the update in
\eqref{eq:raw-compressed-local-update} induces a normalized transition on
realizable states \(K_L\).  If \(K_L=K_{\widetilde L}\), the two raw libraries
have identical project availability, cost, duration, incumbent operating
reward blocks, and joint terminal-belief--next-compressed-state laws after
projection.  Hence \((b,K_L)\) is a controlled semi-Markov projection at
embedded decision epochs and, for every finite calendar horizon \(h\),
\[
  V_h^{\mathrm{raw}}(b,L)=\bar V_h(b,K_L).
\]
Under the additional contraction clause, stationary fixed-point values agree
under the same projection and a maximizing stationary compressed selector
lifts to a Bellman-optimal raw selector.
\end{theorem}

\noindent\emph{Proof sketch.}
Equation~\eqref{eq:raw-compressed-local-update} makes the next compressed state
a deterministic pushforward of the admitted raw update.  The
library-invariance assumptions make this law, every feasible action, and its
payoff block constant on each fiber of \(L\mapsto K_L\).  Strong induction on
calendar horizon is well founded because every research duration is positive;
corresponding raw and compressed action values, and then their finite maxima,
agree.  Bellman intertwining and contraction give the stationary clauses.
The complete normalization and projection proof is in Online Resource~1,
Sections~S1--S2.  The result permits belief--admission dependence and makes no
generic quotient-minimality claim.

\paragraph{Operational factorization audit.}
The assumption can be checked table by table. Project availability, cost,
duration, operation while active, and the joint completion law must be constant
whenever \((b,F_L,C_L)\) is constant; candidate generation and verification
must be constant whenever \((q,b,C_L)\) is constant; and admission must update
the compressed state through \(\operatorname{addK}\). A nontrivial finite class
satisfying these conditions consists of module-gated projects whose
availability is a predicate of \((q,b,C_L)\), whose candidate and verification
tables are indexed by \((q,b,C_L)\), and whose duration, cost, and completion
tables are indexed by \((q,b,F_L,C_L)\). Correlation between terminal belief
and admission remains unrestricted within each indexed joint table.

The condition is substantive. Let two one-active-strategy libraries
\(L_a=\{s_0,a\}\) and \(L_b=\{s_0,b\}\) have identical zero profiles and no
modules. Then \(K_{L_a}=K_{L_b}\). If a zero-cost, one-period project with
deterministic positive completion value is available only when the raw identity
\(a\) is retained, their feasible actions and values differ. The projection
therefore fails. The same failure occurs if a primitive depends directly on
raw cardinality, provenance, owner, or archival status. These variables must
either be included in the compressed state or ruled out by the declared
information architecture.

\subsection{Retention burden and exact source-relative compression}

Let \(w_{s_0}=0\) and give every active strategy an exact rational retention
weight \(w_s>0\).  The one-time burden is
\begin{equation}
  W(L):=\sum_{s\in L}w_s.
  \label{eq:library-burden}
\end{equation}
It may index storage, validation, monitoring, active maintenance, retrieval,
governance, and human attention. Additivity is a modeling assumption. A weight
is not automatically an observed monetary or labor cost, and it is not
operating payoff or project cost. Burden is not a component of \(K_L\): equal
productive states can have different burdens. In the financial cases,
``retained'' means retained in the active maintenance and validation set;
mandatory archival records, code, data, ownership, and reactivation controls
remain outside this binary decision.

For a fixed verified source \(L\), exact source-relative compression chooses
an inactive-containing sublibrary \(L'\subseteq L\) with
\(K_{L'}=K_L\) and minimum \(W(L')\).  By
\cref{thm:raw-to-compressed-projection}, this outer retention decision changes
only the raw representation, not the declared productive process.  Section~\ref{sec:cover-complexity}
defines its feasible family and optimization model.
